\documentclass[11pt,a4paper]{article}
\usepackage[utf8]{inputenc}
\usepackage[T1]{fontenc}
\usepackage[french]{babel}
\usepackage{amsmath,amssymb,amsthm,mathtools,bm}
\usepackage{booktabs}
\usepackage{enumitem}
\usepackage[dvipsnames]{xcolor}
\usepackage[margin=2.5cm]{geometry}
\usepackage[colorlinks=true,linkcolor=blue!60!black,citecolor=green!40!black]{hyperref}

\newtheorem{proposition}{Proposition}
\newtheorem{theoreme}{Théorème}
\newtheorem{lemme}{Lemme}
\theoremstyle{definition}
\newtheorem{hyp}{Hypothèse}
\theoremstyle{remark}
\newtheorem{remarque}{Remarque}

\newcommand{\E}{\mathbb{E}}
\renewcommand{\P}{\mathbb{P}}
\newcommand{\R}{\mathbb{R}}
\newcommand{\dd}{\mathrm{d}}
\newcommand{\sinf}{\sigma_\infty}
\DeclareMathOperator{\erfi}{erfi}
\DeclareMathOperator{\Var}{Var}
\DeclareMathOperator*{\argmax}{arg\,max}

\title{Stratégie de retour à la moyenne sur un processus d'Ornstein--Uhlenbeck :\\
hypothèses, critère ergodique et solutions explicites}
\author{Notes de travail --- projet \texttt{pairs\_trading}}
\date{10 novembre 2026}

\begin{document}
\maketitle

\begin{abstract}
On formalise la stratégie implémentée dans \texttt{research/institutional/}. Un processus scalaire $X$ (le spread)
est supposé suivre un processus d'Ornstein--Uhlenbeck (OU). On considère la famille des politiques de bandes
symétriques $(a,e)$, on évalue leur gain asymptotique par la théorie des processus de renouvellement avec
récompense, et on le maximise. L'optimum est atteint par la politique de \emph{retournement} ($e=a$), pour laquelle
le temps de cycle, le niveau optimal et le gain s'écrivent en fonction de $\erfi$ et d'une intégrale de Dawson.
On explicite ensuite l'estimation des paramètres (biais de Kendall), la déformation introduite par le filtre de
Kalman de couverture, la probabilité de succès par la fonction d'échelle, et un protocole de mesure de la
mauvaise spécification. Chaque résultat est accompagné de son statut : démontré, vérifié numériquement, ou ouvert.
\end{abstract}

\tableofcontents

% =============================================================================
\section{Cadre et hypothèses}
% =============================================================================

Soient $(\Omega,\mathcal F,(\mathcal F_t),\P)$ un espace filtré supportant un mouvement brownien $W$, et
$X=(X_t)_{t\ge 0}$ un processus réel adapté, appelé \emph{spread}. Le temps est mesuré en unités de barre
(pas d'observation $\Delta=1$).

\begin{hyp}[Dynamique OU]\label{h:ou}
$X$ résout $\dd X_t=\theta(\mu-X_t)\,\dd t+\sigma\,\dd W_t$ avec $\theta>0$, $\mu\in\R$, $\sigma>0$ constants.
\end{hyp}

\begin{hyp}[Observation discrète]\label{h:disc}
$X$ n'est observé qu'aux instants entiers, et les décisions sont prises sur ces observations.
\end{hyp}

\begin{hyp}[Gain linéaire]\label{h:pnl}
Une position $q_t\in\{-1,0,+1\}$ portée sur $[t,t+1)$ rapporte $q_t\,(X_{t+1}-X_t)$ par unité de notional.
\end{hyp}

\begin{hyp}[Coût proportionnel]\label{h:cost}
Chaque position (ouverture puis fermeture) coûte $c>0$ par unité de notional, indépendamment de sa durée.
Un retournement $q:\pm1\to\mp1$ ferme une position et en ouvre une.
\end{hyp}

\begin{hyp}[Neutralité au risque]\label{h:rn}
Le critère de choix est l'espérance du gain par unité de temps, sans pénalité de variance.
\end{hyp}

\begin{hyp}[Paramètres connus ou estimés par régression AR(1)]\label{h:est}
Dans le moteur, $(\theta,\mu,\sigma)$ sont réestimés sur une fenêtre glissante de $N$ barres.
\end{hyp}

\begin{remarque}
Les hypothèses \ref{h:ou}, \ref{h:est} et le choix du spread sont testés empiriquement en \S\ref{sec:miss} et
\S\ref{sec:kalman}. L'hypothèse \ref{h:rn} est celle qui produit une politique investie en permanence (\S\ref{sec:opt}).
\end{remarque}

\subsection{Discrétisation exacte}

\begin{proposition}\label{p:ar1}
Sous \ref{h:ou}, avec $b=e^{-\theta}$,
\[
X_{t+1}=\mu+b\,(X_t-\mu)+\varepsilon_{t+1},\qquad
\varepsilon_{t+1}\sim\mathcal N\!\Big(0,\ \tfrac{\sigma^2}{2\theta}(1-b^2)\Big)\ \text{i.i.d.},
\]
et la loi stationnaire est $\mathcal N(\mu,\sinf^2)$ avec $\sinf=\sigma/\sqrt{2\theta}$.
\end{proposition}

L'inversion utilisée pour calibrer est donc
\begin{equation}\label{eq:inversion}
\hat\theta=-\ln\hat b,\qquad \hat\mu=\frac{\hat a}{1-\hat b},\qquad
\hat\sigma=\hat s_\varepsilon\sqrt{\frac{-2\ln\hat b}{1-\hat b^2}},
\end{equation}
où $(\hat a,\hat b)$ sont les coefficients de la régression $X_{t+1}=a+bX_t+\varepsilon$ et $\hat s_\varepsilon$
l'écart-type des résidus.

% =============================================================================
\section{Réduction adimensionnée}
% =============================================================================

Posons $x=(X-\mu)/\sinf$ et $\tilde t=\theta t$. Alors $\tilde x$ résout
\begin{equation}\label{eq:ou-std}
\dd x_{\tilde t}=-x_{\tilde t}\,\dd\tilde t+\sqrt2\,\dd\tilde W_{\tilde t},
\end{equation}
d'espérance nulle et de variance stationnaire $1$. Tout résultat sur \eqref{eq:ou-std} se transporte au cas général
par $a=\sinf\,\tilde a$, $e=\sinf\,\tilde e$, durée $=\tilde T/\theta$ en barres, et le coût devient
\begin{equation}\label{eq:ctilde}
\tilde c=\frac{c}{\sinf}=\frac{c\sqrt{2\theta}}{\sigma}.
\end{equation}
\textbf{Le seul paramètre sans dimension du problème est $\tilde c$.}

\subsection{Générateur, fonction d'échelle, mesure de vitesse}

Le générateur de \eqref{eq:ou-std} est $\mathcal L u=u''-y\,u'$. On pose
\begin{equation}\label{eq:echelle}
s'(y)=e^{y^2/2},\quad S(y)=\int_0^y e^{t^2/2}\,\dd t=\sqrt{\tfrac{\pi}{2}}\ \erfi\!\Big(\tfrac{y}{\sqrt2}\Big),\quad
m(y)=e^{-y^2/2},
\end{equation}
de sorte que $\mathcal L u=s'\,(u'/s')'$. $S$ est strictement croissante, impaire, et $S(y)\sim e^{y^2/2}/y$ en $+\infty$.
Les fonctions harmoniques ($\mathcal L u=0$) sont affines en $S$.

\subsection{Temps de passage}

Soit $\Phi$ la fonction de répartition de $\mathcal N(0,1)$ et $g(y)=e^{y^2/2}\Phi(y)$. On note $T_b$ le temps
d'atteinte de $b$ et $\tau_{(\alpha,\beta)}$ le temps de sortie de $(\alpha,\beta)$.

\begin{proposition}[Passage vers le haut]\label{p:fpt}
Pour $x<b$,
\[
\E_x[T_b]=\sqrt{2\pi}\int_x^{b} g(y)\,\dd y .
\]
\end{proposition}
\begin{proof}[Esquisse]
$u(x)=\E_x[T_b]$ est la solution de $\mathcal L u=-1$, $u(b)=0$, de croissance sous-exponentielle en $-\infty$.
Avec $\mathcal L u=s'(u'/s')'$ on a $(u'/s')'=-e^{-y^2/2}$, d'où $u'(y)/s'(y)=-\int_{-\infty}^{y}e^{-z^2/2}\dd z=-\sqrt{2\pi}\,\Phi(y)$
(la constante est fixée par la condition de croissance). Donc $u'(y)=-\sqrt{2\pi}\,e^{y^2/2}\Phi(y)$ et l'on intègre de $x$ à $b$.
\end{proof}

\begin{proposition}[Sortie d'un intervalle]\label{p:exit}
Pour $-\alpha<x<\beta$, avec $D=S(\beta)-S(-\alpha)$,
\[
\E_x[\tau_{(-\alpha,\beta)}]=\frac{S(\beta)-S(x)}{D}\int_{-\alpha}^{x}\!\big(S(z)-S(-\alpha)\big)m(z)\,\dd z
+\frac{S(x)-S(-\alpha)}{D}\int_{x}^{\beta}\!\big(S(\beta)-S(z)\big)m(z)\,\dd z .
\]
\end{proposition}
(Fonction de Green de $\mathcal L$ avec conditions de Dirichlet; voir Borodin--Salminen, \emph{Handbook of Brownian
Motion}.) Les propositions \ref{p:fpt} et \ref{p:exit} ont été vérifiées contre une résolution de l'EDO par
différences finies (erreur $<10^{-3}$ sur six cas) et par simulation de Monte-Carlo.

\begin{proposition}[Probabilité de sortie par un bord]\label{p:hit}
Pour $-\alpha<x<\beta$,
\[
\P_x\big(T_\beta<T_{-\alpha}\big)=\frac{S(x)-S(-\alpha)}{S(\beta)-S(-\alpha)}.
\]
\end{proposition}
\begin{proof}
$h(x)=\P_x(T_\beta<T_{-\alpha})$ vérifie $\mathcal Lh=0$, $h(-\alpha)=0$, $h(\beta)=1$; les solutions sont affines en $S$.
\end{proof}

% =============================================================================
\section{Critère ergodique et optimum}\label{sec:opt}
% =============================================================================

\subsection{Famille de politiques}

Pour $a>0$ et $e\in(-a,a]$, la politique $\pi_{a,e}$ de \eqref{eq:ou-std} est :
\begin{itemize}[nosep]
\item à plat, on prend une position \emph{longue} dès que $x\le -a$ et \emph{courte} dès que $x\ge a$ ;
\item une position longue est fermée quand $x\ge e$, une position courte quand $x\le -e$ ;
\item si $e=a$, la fermeture d'une position longue en $x=a$ coïncide avec l'ouverture d'une position courte :
c'est le \emph{retournement}, la politique n'est jamais à plat.
\end{itemize}

\subsection{Gain asymptotique par renouvellement}

Par symétrie $x\mapsto -x$, les positions longues et courtes sont équivalentes en loi. Une position longue
entrée en $-a$ et fermée en $e$ rapporte exactement
\begin{equation}\label{eq:profit}
\pi(a,e)=a+e-\tilde c \quad(\text{en unités de }\sinf),
\end{equation}
car $x$ est continu et ne franchit pas de niveau sans l'atteindre. La durée d'un cycle (position + éventuelle attente à plat) est
\begin{equation}\label{eq:cycle}
T(a,e)=\E_{-a}[T_e]+\E_e\big[\tau_{(-a,a)}\big],
\end{equation}
avec la convention que le second terme est nul si $e=a$.

\begin{theoreme}[Renouvellement avec récompense]\label{t:renew}
Les cycles successifs sont i.i.d.\ en loi à une symétrie près et $\E[T]<\infty$. Presque sûrement,
\[
\lim_{H\to\infty}\frac{1}{H}\,(\text{gain cumulé sur }[0,H])=R(a,e):=\frac{a+e-\tilde c}{T(a,e)} .
\]
\end{theoreme}
Il s'agit de la loi des grands nombres pour les processus de renouvellement avec récompense. Le taux
en unités de la barre est $\theta\,\sinf\,R(a,e)$, et les niveaux sont $\mu\pm\sinf\,a$, $\mu\pm\sinf\,e$.

\subsection{Le retournement : forme fermée}

\begin{proposition}\label{p:flip}
Pour la politique de retournement ($e=a$),
\[
T(a,a)=\E_{-a}[T_a]=\sqrt{2\pi}\,S(a)=\pi\,\erfi\!\Big(\tfrac{a}{\sqrt2}\Big).
\]
\end{proposition}
\begin{proof}
Par la proposition \ref{p:fpt}, $T(a,a)=\sqrt{2\pi}\int_{-a}^{a}g$. Or $g(y)+g(-y)=e^{y^2/2}(\Phi(y)+\Phi(-y))=e^{y^2/2}$, donc
$\int_{-a}^{a}g=\int_0^a (g(y)+g(-y))\,\dd y=S(a)$.
\end{proof}

Le gain asymptotique du retournement est donc $R(a)=\dfrac{2a-\tilde c}{\pi\erfi(a/\sqrt2)}$.

\begin{proposition}[Niveau optimal du retournement]\label{p:opt}
Pour tout $\tilde c>0$, $a\mapsto R(a)$ admet un unique maximiseur $a^*(\tilde c)$ sur $(\tilde c/2,\infty)$, caractérisé par
\begin{equation}\label{eq:foc}
\tilde c=\psi(a^*),\qquad \psi(a):=2\Big[a-e^{-a^2/2}S(a)\Big]=2\Big[a-\int_0^a e^{(y^2-a^2)/2}\dd y\Big].
\end{equation}
L'application $\psi$ est une bijection croissante de $(0,\infty)$ sur lui-même, et
$a^*(\tilde c)=\big(\tfrac32\tilde c\big)^{1/3}\,(1+o(1))$ quand $\tilde c\to0$.
\end{proposition}
\begin{proof}
Posons $T(a)=\pi\erfi(a/\sqrt2)$, de sorte que $T'(a)=\sqrt{2\pi}\,e^{a^2/2}$ et $T''>0$.
Alors $R'(a)=h(a)/T(a)^2$ avec $h(a)=2T(a)-(2a-\tilde c)\,T'(a)$. On calcule
$h'(a)=-(2a-\tilde c)\,T''(a)<0$ pour $2a>\tilde c$, $h(\tilde c/2)=2T(\tilde c/2)>0$ et $h(a)\to-\infty$
(car $T'/T\to\infty$ à l'infini). Donc $h$ s'annule une seule fois, $R$ croît puis décroît.
L'équation $h(a)=0$ s'écrit $2\sqrt{2\pi}\,S(a)=(2a-\tilde c)\sqrt{2\pi}\,e^{a^2/2}$, soit \eqref{eq:foc}.
Enfin $\psi'(a)=2a\,e^{-a^2/2}S(a)>0$, $\psi(0)=0$, et le développement
$e^{-a^2/2}S(a)=a-a^3/3+O(a^5)$ donne $\psi(a)=\tfrac23a^3+O(a^5)$.
\end{proof}

En unités dimensionnelles, pour $\tilde c\ll1$ :
\begin{equation}\label{eq:scaling}
a^*\simeq\Big(\frac{3c}{2}\Big)^{1/3}\sigma^{2/3}(2\theta)^{-1/3}\propto c^{1/3}\,\sigma^{2/3}\,\theta^{-1/3},
\end{equation}
la loi d'échelle de Bertram : la bande croît avec le coût et la volatilité, décroît avec la vitesse de rappel.

\begin{table}[h]
\centering
\caption{Retournement optimal (forme fermée, vérifiée contre l'optimisation numérique sur $(a,e)$).}
\begin{tabular}{rrrrr}
\toprule
$\tilde c$ & $a^*$ & $(3\tilde c/2)^{1/3}$ & $T(a^*)$ & $R(a^*)$\\
\midrule
0{,}02 & 0{,}313 & 0{,}311 & 0{,}797 & 0{,}760\\
0{,}10 & 0{,}542 & 0{,}531 & 1{,}427 & 0{,}689\\
0{,}30 & 0{,}799 & 0{,}766 & 2{,}237 & 0{,}580\\
1{,}00 & 1{,}264 & 1{,}145 & 4{,}262 & 0{,}359\\
\bottomrule
\end{tabular}
\end{table}

\subsection{Optimisation sur $(a,e)$ : statut}

Le taux $R(a,e)$ des équations \eqref{eq:profit}--\eqref{eq:cycle} a été calculé sur une grille
(pas $0{,}01$ en $a$ et $e$, $a\le6$) à l'aide des propositions \ref{p:fpt} et \ref{p:exit} par cumuls trapézoïdaux.

\begin{remarque}[Résultat numérique, non démontré]
Pour les six valeurs $\tilde c\in\{0{,}02;\,0{,}1;\,0{,}3;\,0{,}6;\,1;\,1{,}5\}$, le maximum sur $(a,e)$ est atteint en $e=a$.
La politique « sortie à la moyenne » ($e=0$) atteint $97\,\%$, $91\,\%$, $81\,\%$, $67\,\%$, $51\,\%$, $33\,\%$ du taux optimal.
Une preuve de $\arg\max_e R(a,e)=a$ pour tout $a$ reste ouverte.
\end{remarque}

La théorie a été confrontée à une simulation directe de la politique sur un OU exact (pas $0{,}005$, $6\cdot10^6$ pas) :
$R_{\mathrm{th}}=0{,}689$ contre $0{,}694$ pour $\tilde c=0{,}1$, et $0{,}468$ contre $0{,}465$ pour $\tilde c=0{,}6$.

% =============================================================================
\section{Critère actualisé (Leung--Li)}
% =============================================================================

Pour un taux $r>0$, le problème d'arrêt optimal de Leung et Li (2015) pour un OU avec coût $c$ se résout en deux étapes.

\paragraph{Sortie d'une position longue.} On cherche $V(x)=\sup_{\tau}\E_x\big[e^{-r\tau}(X_\tau-c)\big]$.
Soient $F$ et $G$ les solutions croissante et décroissante de $\tfrac{\sigma^2}{2}f''+\theta(\mu-x)f'-rf=0$ ; en termes de la
fonction cylindre parabolique $D_\nu$ de Weber, avec $z=(x-\mu)\sqrt{2\theta}/\sigma$ et $\nu=-r/\theta$,
\[
F(x)=e^{z^2/4}D_{\nu}(-z),\qquad G(x)=e^{z^2/4}D_{\nu}(z).
\]
Alors $V(x)=(b^*-c)\,F(x)/F(b^*)$ sur $x<b^*$, où $b^*=\argmax_b\,(b-c)/F(b)$.

\paragraph{Entrée.} $J(x)=\sup_{\nu}\E_x\big[e^{-r\nu}(V(X_\nu)-X_\nu-c)\big]$ s'écrit $J(x)=\big(V(d^*)-d^*-c\big)\,G(x)/G(d^*)$ pour $x>d^*$, avec
$d^*=\argmax_d\,\big(V(d)-d-c\big)/G(d)$.

\begin{remarque}[Dégénérescence]
Les résidus de l'EDO pour $F,G$ ont été vérifiés ($<4\cdot10^{-7}$ en relatif). Pour $\theta=0{,}1$, $\sigma=0{,}01$, $c=10^{-3}$ :
$r/\theta=10^{-1},\,10^{-2},\,10^{-3},\,5{,}7\cdot10^{-5}$ donnent des niveaux d'entrée de $1{,}36,\,2{,}20,\,2{,}94,\,3{,}72$ écarts-types.
Pour un taux annuel de $5\,\%$ converti par barre horaire, $r/\theta\sim10^{-4}$ et les niveaux partent vers la borne : « attendre est gratuit ».
Les niveaux $b^*,d^*$ sont des niveaux absolus de $X$ et le problème n'est pas invariant par translation quand $r>0$ (le gain $X_\tau-c$ est un capital détenu).
Le critère actualisé n'est donc pas adapté à l'horizon intra-journalier avec un taux réaliste. C'est la raison pour laquelle le critère ergodique de la \S\ref{sec:opt}
est le critère par défaut.
\end{remarque}

% =============================================================================
\section{Estimation des paramètres}
% =============================================================================

\subsection{Biais de Kendall}

\begin{proposition}[Kendall, 1954]
Pour une régression AR(1) avec constante estimée sur $N$ observations,
$\E[\hat b]=b-\dfrac{1+3b}{N}+O(N^{-2})$.
\end{proposition}

Pour $\theta=0{,}02$ ($b=0{,}980$) et $N=156$ : $b-(1+3b)/N=0{,}955$, d'où $\hat\theta\approx0{,}046$,
contre une médiane de $0{,}047$ mesurée sur $4000$ simulations. Le biais va dans le sens d'un rappel \emph{plus rapide}
que le vrai. Comme $\tilde c=c\sqrt{2\theta}/\sigma$ dépend de $\hat\theta$ et $\hat\sigma$ (équation \eqref{eq:inversion}) et que les niveaux sont
$a^*(\tilde c)\,\sigma/\sqrt{2\theta}$, l'effet net sur les seuils n'est pas dérivé ici ; il est mesuré indirectement en \S\ref{sec:miss} (écart moteur--oracle).

\subsection{Incertitude de $\hat\theta$}

L'erreur-type asymptotique de $\hat b$ est $\sqrt{(1-b^2)/N}$. Pour $\theta\approx0{,}04$, $N=156$ elle vaut environ $0{,}023$ sur $\theta$ :
la dispersion de $\hat\theta$ d'une fenêtre à l'autre contient donc une part de bruit d'échantillonnage du même ordre que la dispersion observée
(rapport des déciles $9\times$ sur ADA/AVAX).

% =============================================================================
\section{Couche de couverture par filtre de Kalman}\label{sec:kalman}
% =============================================================================

Le spread négocié n'est pas $X$ mais une \emph{innovation} de filtre. Avec $y_t=\log P^A_t$, $h_t=(\log P^B_t,\,1)^\top$ et un état
$\vartheta_t=(\beta_t,\alpha_t)^\top$ évoluant en marche aléatoire de covariance $Q$, le filtre calcule
\[
e_t=y_t-h_t^\top\hat\vartheta_{t|t-1},\qquad S_t=h_t^\top P_{t|t-1}h_t+R,\qquad
\hat\vartheta_{t|t}=\hat\vartheta_{t|t-1}+K_t e_t .
\]

\begin{proposition}[Effet passe-haut, approximation scalaire]
Si $\alpha$ est estimé avec un gain constant $K\in(0,1)$ et si $\beta$ est supposé exact, alors $e_t=X_t-m_{t-1}$ avec
$m_t=(1-K)m_{t-1}+KX_t$, c'est-à-dire
\[
e=\frac{1-L}{1-(1-K)L}\,X ,
\]
où $L$ est l'opérateur retard. La fonction de transfert s'annule en fréquence nulle : les composantes lentes de $X$ sont supprimées.
\end{proposition}
\begin{proof}
$m_{t-1}=K\sum_{j\ge0}(1-K)^jX_{t-1-j}$, d'où $e=X-\frac{KL}{1-(1-K)L}X=\frac{1-(1-K)L-KL}{1-(1-K)L}X$.
\end{proof}

Conséquence : l'autocorrélation d'ordre 1 de l'innovation est inférieure à celle de $X$. Pour $X$ AR(1) de coefficient $\varphi=0{,}9$ :

\begin{center}
\begin{tabular}{lcccc}
\toprule
$K$ & $0$ & $0{,}05$ & $0{,}10$ & $0{,}20$\\
$\rho_e(1)$ & $0{,}900$ & $0{,}851$ & $0{,}804$ & $0{,}708$\\
\bottomrule
\end{tabular}
\end{center}

Sur une paire simulée avec $\varphi=0{,}9$ le code mesure $\hat\varphi_e=0{,}736$, ce qui correspond à $K_{\mathrm{eq}}\approx0{,}17$ dans ce modèle
(approximation : l'estimation conjointe de $\beta$ n'est pas prise en compte).

\paragraph{Pourquoi c'est un problème d'estimation et de décision.} Les seuils \eqref{eq:foc} sont calculés avec les paramètres de $e$ (donc $\hat\theta_e>\theta$, $\hat\sigma_e<\sigma$)
alors que le gain réalisé est $\Delta X$ (hypothèse \ref{h:pnl}) : le signal est calculé sur une série dont la dynamique n'est pas celle qu'on encaisse.

% =============================================================================
\section{Dimensionnement : probabilité de succès et critère de Kelly}
% =============================================================================

Dans \eqref{eq:ou-std} un trade gagne \emph{exactement} $a+e-\tilde c$ : il n'y a pas de perte intrinsèque.
On introduit une perte adverse hypothétique : le trade est abandonné si $x$ atteint $-(a+s)$, $s>0$ (en unités de $\sinf$), avec perte $s+\tilde c$.

\begin{proposition}
Par la proposition \ref{p:hit}, la probabilité de succès est
\[
p(a,e,s)=\frac{S(-a)-S(-a-s)}{S(e)-S(-a-s)} .
\]
\end{proposition}
Vérifiée par Monte-Carlo (écart $<0{,}03$ avec $4000$ trajectoires). Avec $w=a+e-\tilde c$, $\ell=s+\tilde c$, $\beta=w/\ell$, la fraction de Kelly est
\[
f^*=\frac{p\,\beta-(1-p)}{\beta},\qquad f=\min\big(\lambda f^*_+,\ f_{\max}\big),\quad \lambda\in(0,1].
\]
\begin{remarque}
$f^*$ dépend fortement de $s$, qui est une hypothèse de risque et non un résultat du modèle : sur une paire simulée,
$\lambda f^*=0{,}04;\ 0{,}20;\ 0{,}25$ pour $s=1;\ 2;\ 3$ écarts-types ($\lambda=0{,}25$).
\end{remarque}

% =============================================================================
\section{Mauvaise spécification : protocole et premiers résultats}\label{sec:miss}
% =============================================================================

\subsection{Processus générateurs}

Les écarts à \ref{h:ou} sont introduits un par un puis ensemble, avec $x_{t+1}=\mu_{t+1}+b_{t+1}(x_t-\mu_t)+\sigma_{t+1}\xi_{t+1}+J_{t+1}$ et $b_t=e^{-\theta_t}$ :
\begin{itemize}[nosep]
\item queues : $\xi$ de loi de Student à $\nu$ degrés de liberté, normalisée en variance $1$ ;
\item volatilité stochastique : $\ln\sigma_t=\ln\sigma_0+h_t-\tfrac12\Var h$, $h_{t+1}=\rho h_t+\eta\,\zeta_{t+1}$ ;
\item sauts : $J_t=\mathbf 1_{\{U_t<p_J\}}\,\kappa\,\sigma_0\,\zeta'_t$ ;
\item niveau moyen dérivant : $\mu_{t+1}=\mu_t+\delta\,\sigma_0\,\zeta''_{t+1}$ ;
\item régimes : $\theta_t\in\{\theta_{\mathrm{h}},\theta_{\mathrm{b}}\}$, chaîne de Markov à deux états.
\end{itemize}
La calibration vise les statistiques mesurées sur le spread réel ADA/AVAX :

\begin{center}
\begin{tabular}{lrr}
\toprule
 & réel & scénario combiné\\
\midrule
$\hat\theta$ médian ($N=156$) & $0{,}037$ & $0{,}050$\\
rapport des déciles de $\hat\theta$ & $9{,}0$ & $7{,}2$\\
kurtosis excédentaire des résidus (médiane) & $2{,}6$ & $3{,}9$\\
autocorrélation de $\hat\varepsilon^2$ (médiane) & $0{,}10$ & $0{,}09$\\
\bottomrule
\end{tabular}
\end{center}

\subsection{Références de comparaison}

\begin{itemize}[nosep]
\item \textbf{Moteur} : calibration AR(1) glissante \eqref{eq:inversion}, bandes optimales $\pm a^*\sinf$.
\item \textbf{Oracle} : mêmes bandes avec les paramètres \emph{vrais} $(\theta_t,\mu_t,\sigma_t)$ à chaque barre ; mesure le coût de l'estimation et de la dérive.
\item \textbf{Échelle $k$} : bandes multipliées par $k$ ; si $\arg\max_k=1$ la forme OU des seuils est adaptée.
\end{itemize}

\subsection{Deux défauts de l'implémentation, isolés par ce protocole}

Les premiers essais montraient un moteur perdant même avec les \emph{vrais} paramètres. L'expérience contrôlée a isolé deux causes, toutes deux étrangères à la forme OU :

\begin{enumerate}[nosep]
\item \textbf{Référence des seuils.} Les niveaux $\mu\pm a\sinf$ étaient placés autour d'une moyenne exponentielle rapide $m_t$ du spread, et non autour du $\mu$ de la calibration. Le signal effectif était
$x_t-m_t$, c'est-à-dire de nouveau un filtre passe-haut (même mécanisme que la proposition ci-dessus), avec les gains correspondants :
à paramètres vrais et sans couche de Kalman, rendement $84\,\%$ (seuils sur $m_t$) contre $337\,\%$ (seuils sur $\mu$ calibré), la théorie prédisant $335\,\%$
($147\,\%$ additifs capitalisés sur $3400$ barres, $\tilde c=0{,}066$). Avec $\mu$ calibré, le facteur d'échelle optimal des bandes est $k^*\approx1$ (plat entre $0{,}7$ et $1$), ce qui lève l'écart noté plus haut.
\item \textbf{Couche de couverture.} L'innovation de Kalman est remplacée par $x_t=\log P^A_t-\hat\beta\log P^B_t-\hat\alpha$ où $(\hat\beta,\hat\alpha)$ est la régression par moindres carrés avec constante sur les $500$ barres
strictement passées, rafraîchie toutes les $24$ barres. L'autocorrélation d'ordre 1 du spread est préservée ($\hat\varphi=0{,}979$ pour $0{,}980$) au lieu de $0{,}72$.
La constante évite un saut de niveau $\Delta\beta\,\log P^B$ aux rafraîchissements.
\end{enumerate}

Deux réglages ont de plus été mesurés. \emph{Fenêtre de calibration} $N$ : sur un OU pur (arrêt temporel $2\times$ demi-vie) le rendement moyen passe de $197\,\%$ ($N=156$) à $435\,\%$ ($N=312$) puis $487\,\%$ ($N=624$), l'oracle étant à $636\,\%$ ;
l'ordre est le même avec dérive du niveau moyen et dans le scénario combiné. Justification : pour $\theta=0{,}02$, une fenêtre de $156$ barres contient environ $N\theta/2\approx1{,}6$ échantillons indépendants, l'écart-type de $\hat\mu$ est de l'ordre de $\sinf$ alors que la demi-largeur de bande optimale vaut $a^*\sinf\approx0{,}47\sinf$.
\emph{Arrêt temporel} : $2$, $4$, $8$ demi-vies ou aucun donnent $345$, $364$, $381$, $386\,\%$ sur OU pur, et $109$, $117$, $119$, $119\,\%$ avec une rupture de cointégration à $60\,\%$ de la série (pire drawdown identique, $-38\,\%$) ; défaut retenu : $8$ demi-vies.

\subsection{Résultats avec les réglages corrigés}

Rendement total moyen en \% sur $16$ trajectoires de $6000$ barres ($\beta=1{,}5$ inconnu du moteur, position $=100\,\%$ du capital, $c=20$ pb).
« Historique » : couverture de Kalman, seuils sur la moyenne rapide, $N=156$, arrêt à $2$ demi-vies. « Nouveau » : couverture OLS roulante, $\mu$ calibré, $N=624$, arrêt à $8$ demi-vies.
« Oracle » : $\beta$ et $(\theta_t,\mu_t,\sigma_t)$ vrais à chaque barre.

\begin{center}
\small
\begin{tabular}{lrrrrrr}
\toprule
scénario & historique & nouveau & $k=0{,}7$ & $k=1{,}5$ & oracle & nouveau/oracle\\
\midrule
OU pur & $-64$ & $381$ & $384$ & $379$ & $715$ & $53\,\%$\\
+ queues de Student & $-48$ & $391$ & $389$ & $376$ & $701$ & $56\,\%$\\
+ volatilité stochastique & $-54$ & $345$ & $336$ & $359$ & $1176$ & $29\,\%$\\
+ sauts & $-66$ & $487$ & $436$ & $445$ & $939$ & $52\,\%$\\
+ niveau moyen dérivant & $-63$ & $401$ & $359$ & $365$ & $672$ & $60\,\%$\\
+ régimes de $\theta$ & $-60$ & $316$ & $353$ & $276$ & $865$ & $36\,\%$\\
rupture de cointégration & $-66$ & $119$ & $124$ & $144$ & $179$ & $66\,\%$\\
tout combiné & $-42$ & $464$ & $419$ & $480$ & $1636$ & $28\,\%$\\
\bottomrule
\end{tabular}
\end{center}

\paragraph{Lecture.}
\begin{enumerate}[nosep]
\item La configuration historique perd dans tous les scénarios, y compris l'OU pur ; la nouvelle gagne dans tous (en moyenne), sur $88$ à $100\,\%$ des trajectoires.
\item Les seuils théoriques ($k=1$) ne sont pas nettement battus par $k=0{,}7$ ou $k=1{,}5$ : les différences ne sont pas systématiques d'un scénario à l'autre (aucune erreur-type n'a été calculée ici ; avec $16$ trajectoires de rendement très dispersé, elles ne sont pas établies).
\item Le coût résiduel de l'estimation est important : le moteur capte $28$ à $66\,\%$ de l'oracle. Il est le plus fort là où les paramètres varient vite ($\sigma_t$ stochastique, régimes de $\theta$, scénario combiné : $28$ à $36\,\%$),
ce qui est cohérent avec le compromis d'une longue fenêtre (précision) contre sa lenteur d'adaptation. Une estimation adaptative est donc la piste suivante.
\item L'oracle de la volatilité stochastique connaît $\sigma_t$ à chaque barre : il est plus favorable que ce qu'un estimateur causal peut atteindre ; le rapport $29\,\%$ est un minorant du rapport atteignable.
\end{enumerate}

% =============================================================================
\section{Récapitulatif : ce qui est établi et ce qui est ouvert}
% =============================================================================

\begin{center}
\begin{tabular}{p{0.52\textwidth}p{0.38\textwidth}}
\toprule
Énoncé & Statut\\
\midrule
Discrétisation exacte, inversion \eqref{eq:inversion} & démontré\\
Temps de passage, probabilité de sortie (prop.\ \ref{p:fpt}--\ref{p:hit}) & démontré (esquisse) + vérifié numériquement\\
Gain asymptotique par renouvellement (th.\ \ref{t:renew}) & résultat classique\\
Cycle du retournement $=\pi\erfi(a/\sqrt2)$, optimum unique \eqref{eq:foc} & démontré\\
Optimalité de $e=a$ sur la famille $(a,e)$ & vérifié sur grille, preuve ouverte\\
Biais de Kendall sur $\hat\theta$ & résultat classique, vérifié par simulation\\
Effet passe-haut de l'innovation de Kalman & démontré (approximation scalaire), effet mesuré\\
Origine de $k^*\neq1$ à paramètres connus & expliqué : seuils ancrés sur une moyenne rapide ; $k^*\approx1$ après correction\\
Estimation adaptative de $(\theta,\sigma,\mu)$ (volatilité stochastique, régimes) & \textbf{ouvert}\\
Politique optimale hors OU (queues, volatilité stochastique, régimes) & \textbf{ouvert}\\
\bottomrule
\end{tabular}
\end{center}

\begin{thebibliography}{9}
\bibitem{bertram} W.\ K.\ Bertram, \emph{Analytic solutions for optimal statistical arbitrage trading}, Physica A 389 (2010) 2234--2243.
\bibitem{leungli} T.\ Leung, X.\ Li, \emph{Optimal mean reversion trading with transaction costs and stop-loss exit}, Int.\ J.\ Theor.\ Appl.\ Finance 18 (2015).
\bibitem{kendall} M.\ G.\ Kendall, \emph{Note on bias in the estimation of autocorrelation}, Biometrika 41 (1954) 403--404.
\bibitem{borodin} A.\ N.\ Borodin, P.\ Salminen, \emph{Handbook of Brownian Motion --- Facts and Formulae}, Birkhäuser, 2002.
\bibitem{kelly} J.\ L.\ Kelly, \emph{A new interpretation of information rate}, Bell Syst.\ Tech.\ J.\ 35 (1956).
\end{thebibliography}

\end{document}
